\chapter{Syncope}

\section*{Definition}
Transient loss consciousness temporary reduction cerebral perfusion spontaneous recovery.

\section*{Pathophysiology}
Differentiate vasovagal cardiac neurologic metabolic causes syncope exertion particularly concerning requires cardiac evaluation.

\section*{Etiology}
\begin{itemize}
  \item Vasovagal episode
  \item Cardiac arrhythmia
  \item Orthostatic hypotension
  \item Carotid sinus hypersensitivity
  \item Seizure mimic
\end{itemize}

\section*{Indications for Evaluation}
Seek care if:
\begin{itemize}
  \item Severe or worsening symptoms
  \item Symptoms lasting more than First episode or recurrent with injury risk warrants medical evaluation especially if exertional precipitated by position changes cardiac family history present
  \item Accompanied by other concerning signs
\end{itemize}

\section*{Related Symptoms}
\begin{itemize}
  \item Palpitations
  \item Chest pain
  \item Shortness of breath
  \item Postural dizziness
  \item Prodromal warning symptoms
\end{itemize}
\textbf{Note:} Always consider serious underlying conditions when evaluating persistent syncope.

\section*{Self-Care / Home Management}
\begin{itemize}
  \item Ensure a safe environment to prevent injury during episodes (remove hazards, avoid driving or operating machinery).
  \item Keep a symptom diary documenting triggers, duration, frequency, and associated features.
  \item Practice stress reduction techniques (deep breathing, meditation, adequate sleep) as stress often exacerbates neurological symptoms.
  \item Avoid alcohol, caffeine, and recreational drugs which can worsen many neurological conditions.
  \item Seek immediate evaluation for sudden-onset, severe, or progressive neurological symptoms.
\end{itemize}

\section*{Diagnostic Approach}
\begin{itemize}
  \item Detailed neurological history includes onset pattern, progression, triggers, associated symptoms, and functional impact.
  \item Comprehensive neurological examination assesses mental status, cranial nerves, motor/sensory function, reflexes, coordination, and gait.
  \item Neuroimaging (CT, MRI) is often indicated to evaluate for structural, vascular, or demyelinating causes.
  \item Specialized testing (EEG for seizures, nerve conduction studies for neuropathy, lumbar puncture for meningitis) is guided by clinical suspicion.
\end{itemize}

\section*{Treatment / Management Overview}
\begin{itemize}
  \item Treatment is directed at the underlying cause identified through diagnostic evaluation.
  \item Symptomatic pharmacotherapy may include analgesics, anticonvulsants, antidepressants, or disease-modifying agents as appropriate.
  \item Physical, occupational, and speech therapy play important roles in functional recovery and rehabilitation.
  \item Referral to neurology is indicated for unexplained, progressive, or treatment-refractory neurological symptoms.
\end{itemize}

\clearpage
